\section{연구 방법}

% --------------------------------------------------------------------
% 이 파일의 역할
% --------------------------------------------------------------------
% 3-Methods.tex에는 연구가 어떻게 수행되었는지를 작성합니다.
% 연구 설계, 연구 대상, 연구 절차, 자료 수집 방법, 자료 분석 방법,
% 타당도와 신뢰도 확보 방법, 연구 윤리 등을 구체적으로 설명합니다.
% 독자가 같은 절차로 연구를 이해하거나 반복할 수 있을 정도로 명확하게 씁니다.
%
% --------------------------------------------------------------------
% 연구 방법 장 작성 안내
% --------------------------------------------------------------------
% 이 장에서는 연구자가 어떤 절차와 방법으로 연구 문제에 답했는지를
% 다른 사람이 재현하거나 평가할 수 있을 정도로 구체적으로 설명합니다.
% 단순히 "무엇을 했다"를 나열하기보다, 왜 그 방법이 연구 목적에 적절한지
% 함께 밝혀 주는 것이 좋습니다.
%
% 보통 아래 순서로 작성합니다.
% 1. 연구 설계: 연구 유형, 전체 흐름, 연구 문제와 방법의 연결
% 2. 연구 대상: 참여자/자료/문헌/사례의 선정 기준과 특성
% 3. 연구 절차: 연구가 진행된 단계와 일정
% 4. 자료 수집: 어떤 자료를 어떤 도구와 방식으로 수집했는지
% 5. 자료 분석: 수집한 자료를 어떤 기준, 방법, 통계, 코딩 절차로 분석했는지
% 6. 타당도와 신뢰도: 도구 검증, 전문가 검토, 삼각검증, 평정자 일치도 등
% 7. 연구 윤리: 동의, 익명화, 개인정보 보호, IRB 여부 등
%
% 양적 연구라면 표본, 검사 도구, 변인, 통계 분석 방법을 구체적으로 씁니다.
% 질적 연구라면 참여자 선정, 자료 수집 맥락, 코딩 절차, 해석의 신뢰성 확보
% 방법을 자세히 씁니다. 혼합 연구라면 양적/질적 자료가 어떤 순서와 방식으로
% 결합되는지 설명합니다.

연구 방법 장은 과학교육 연구자가 어떤 절차와 방법으로 과학 교수·학습 현상에 대한 연구 문제에 답했는지를 설명하는 곳이다.
독자가 연구의 타당성을 평가할 수 있도록 연구 설계, 참여자와 수업 맥락, 절차, 자료 수집, 자료 분석, 타당도와 신뢰도, 연구 윤리를 구체적으로 쓴다.
단순히 수행한 일을 나열하기보다 그 방법이 학습자의 개념 이해, 탐구 수행, 수업 상호작용, 교사 인식처럼 연구 목적에서 다루는 과학교육 현상을 설명하는 데 왜 적절한지도 함께 밝혀야 한다.

이 템플릿에서는 연구 방법을 \verb|sub/3-Methods.tex|에 작성한다.
연구 문제와 방법의 대응 관계는 표로 정리할 수 있으며, 표 번호와 본문 참조는 \verb|\caption{}|, \verb|\label{}|, \verb|\ref{}|로 자동 관리한다.
표의 폭과 열 너비는 \verb|tabular*|, \verb|p{...}|, \verb|\linewidth|를 이용해 조정한다.

\subsection{연구 설계}

% 연구의 전체 설계를 설명합니다.
% 예: 본 연구는 설문 조사 연구, 실험 연구, 개발 연구, 사례 연구,
% 델파이 연구, 혼합 연구 중 어떤 성격을 가지는지 밝힙니다.
% 연구 문제별로 어떤 자료와 분석 방법을 사용했는지도 함께 연결해 줍니다.

연구 설계 절에서는 연구가 어떤 유형의 과학교육 연구인지 먼저 밝힌다.
예를 들어 과학 수업 처치의 효과를 살펴보는 실험 연구, 학생의 인식이나 태도를 조사하는 조사 연구, 특정 수업이나 교사 공동체를 깊이 살펴보는 사례 연구, 교수·학습 자료나 평가 도구를 만드는 개발 연구, 선행 연구 동향을 분석하는 문헌 연구, 전문가 합의를 구하는 델파이 연구, 양적·질적 자료를 함께 다루는 혼합 연구 중 어디에 해당하는지 설명한다.
그다음 연구 목적과 연구 문제에 비추어 왜 해당 연구 방법이 적절한지 논리적으로 제시한다.

\subsubsection{과학교육 연구방법론의 주요 유형}

과학교육 분야에서는 연구 문제가 과학 개념 이해의 변화, 탐구 수행, 과학 학습 동기, 수업 상호작용, 교사의 전문성, 교육과정과 평가의 실행처럼 다양한 층위를 가진다.
따라서 연구 방법은 단순히 선호에 따라 고르는 것이 아니라, 연구 문제가 요구하는 증거의 성격에 맞게 선택해야 한다.
집단 간 차이, 처치 효과, 변인 간 관계를 설명하려면 양적 연구가 적절하고, 학습자의 사고 과정이나 수업 장면의 의미를 깊이 해석하려면 질적 연구가 적절하다.
두 접근이 서로 보완되어야 할 때는 혼합 연구를 설계할 수 있다.

\begin{table}[htbp]
    \centering
    \caption{과학교육 연구방법론의 주요 유형}
    \label{tab:science-education-method-types}
    \ThesisTableStyle
    \begin{tabularx}{\linewidth}{@{}L{0.16\linewidth}Y Y@{}}
        \toprule
        연구 유형 & 주로 다루는 질문 & 과학교육 연구의 자료 예시 \\
        \midrule
        양적 연구 & 수업 처치의 효과, 집단 간 차이, 변인 간 관계를 수치로 설명하는 연구 & 개념 검사, 성취도 검사, 탐구 능력 검사, 과학 태도 설문, 학습 동기 척도 \\
        \midrule
        질적 연구 & 학습자의 사고, 수업 상호작용, 교사의 인식, 특정 사례의 맥락과 의미를 해석하는 연구 & 면담 전사본, 수업 관찰 기록, 활동지, 보고서, 담화 자료, 연구자 관찰 일지 \\
        \midrule
        혼합 연구 & 양적 결과와 질적 해석을 함께 사용하여 연구 문제를 더 입체적으로 설명하는 연구 & 사전·사후 검사와 면담, 설문 결과와 관찰 자료, 학습 산출물 분석과 통계 결과 \\
        \midrule
        개발 연구 & 교수·학습 자료, 평가 도구, 수업 모형, 프로그램을 개발하고 타당화하는 연구 & 전문가 검토 자료, 예비 적용 결과, 수정 이력, 수업 적용 자료, 만족도 조사 \\
        \midrule
        문헌 연구 & 선행 연구의 흐름, 쟁점, 연구 동향, 개념적 틀을 종합하는 연구 & 학술지 논문, 학위논문, 교육과정 문서, 교과서, 연구 보고서 \\
        \midrule
        델파이 연구 & 전문가 의견을 반복적으로 수렴하여 기준, 요소, 역량, 문항의 합의를 도출하는 연구 & 전문가 설문, 의견 수렴지, 내용 타당도 지수, 합의도와 수렴도 지표 \\
        \bottomrule
    \end{tabularx}
\end{table}

표~\ref{tab:science-education-method-types}은 과학교육 연구에서 자주 사용하는 연구방법론을 정리한 것이다.
실제 논문에서는 자신의 연구가 어느 한 유형에만 완전히 들어맞는다고 단정하기보다, 연구 문제와 자료의 성격에 따라 어떤 접근을 주된 방법으로 삼고 어떤 방법을 보조적으로 사용하는지 명확히 밝히는 것이 좋다.
예를 들어 과학 탐구 수업 프로그램의 효과를 확인하는 연구라면 사전·사후 검사 점수를 비교하는 양적 분석을 중심으로 하되, 학생 면담이나 활동지 분석을 통해 점수 변화의 의미를 해석할 수 있다.
반대로 특정 교사의 과학 수업 실행을 깊이 이해하려는 연구라면 질적 사례 연구를 중심으로 하되, 설문이나 검사 결과를 보조 자료로 활용할 수 있다.

연구 설계를 쓸 때는 연구 문제와 방법이 서로 대응되도록 작성하는 것이 중요하다.
아래 표처럼 연구 문제별로 자료 수집 방법과 분석 방법을 정리하면 독자가 연구의 전체 구조를 쉽게 이해할 수 있다.

\begin{table}[htbp]
    \centering
    \caption{연구 문제와 연구 방법의 대응 예시}
    \label{tab:method-design-example}
    \ThesisTableStyle
    \begin{tabular*}{\linewidth}{@{\extracolsep{\fill}}p{0.27\linewidth}p{0.25\linewidth}p{0.30\linewidth}@{}}
        \toprule
        연구 문제 & 자료 수집 방법 & 자료 분석 방법 \\
        \midrule
        연구 문제 1 & 문헌 자료, 선행 연구 & 내용 분석, 범주화 \\
        \midrule
        연구 문제 2 & 설문 조사, 검사 도구 & 기술통계, 평균 비교 \\
        \midrule
        연구 문제 3 & 면담, 관찰 기록 & 전사, 코딩, 주제 분석 \\
        \bottomrule
    \end{tabular*}
\end{table}

표~\ref{tab:method-design-example}은 연구 문제와 연구 방법을 연결하여 제시하는 방식의 예시이다.
실제 논문에서는 자신의 연구 문제에 맞게 자료의 종류, 수집 절차, 분석 방법을 구체적으로 수정한다.

\subsection{연구 대상}

% 연구 참여자, 분석 자료, 문헌, 수업 사례 등의 선정 기준과 특성을 씁니다.
% 사람을 대상으로 한 연구라면 학교급, 학년, 인원, 성별, 모집 방법,
% 참여 조건, 제외 기준 등을 제시합니다.
% 문헌 연구라면 검색 데이터베이스, 검색어, 포함/제외 기준, 최종 선정 편수를 씁니다.

연구 대상 절에서는 누구 또는 무엇을 대상으로 과학교육 연구를 수행했는지 설명한다.
사람을 대상으로 한 연구라면 참여자의 학교급, 학년, 과목, 수업 단원, 인원, 모집 방법, 참여 조건, 제외 기준을 제시한다.
문헌이나 자료를 대상으로 한 연구라면 과학교육 관련 자료를 검색한 데이터베이스, 검색어, 기간, 포함 기준과 제외 기준을 밝힌다.

연구 대상은 단순히 인원수만 제시하는 데 그치지 않는다.
왜 그 대상이 연구 문제에 적절한지, 어떤 기준으로 선정했는지, 연구 결과를 해석할 때 고려해야 할 대상의 특성은 무엇인지 함께 설명한다.
예를 들어 특정 지역, 특정 학년, 특정 수업 맥락에서 수집한 자료라면 그 범위가 연구 결과의 해석과 일반화에 어떤 영향을 주는지도 간단히 밝힌다.

\subsection{연구 절차}

% 연구가 어떤 단계로 진행되었는지 시간 순서대로 설명합니다.
% 필요하면 연구 절차를 표나 그림으로 제시하면 독자가 흐름을 이해하기 쉽습니다.
% 예: 문헌 검토 -> 도구 개발 -> 전문가 검토 -> 예비 조사 -> 본 조사 -> 자료 분석

연구 절차 절에서는 연구가 진행된 과정을 시간 순서대로 설명한다.
보통 문헌 검토, 연구 도구 개발, 전문가 검토, 예비 조사, 본 조사, 자료 분석, 결과 해석의 순서로 정리할 수 있다.
각 단계에서는 무엇을 했는지뿐 아니라 그 단계가 다음 단계와 어떻게 연결되는지도 밝혀야 한다.

예를 들어 검사 도구를 개발한 연구라면 먼저 선행 연구를 바탕으로 문항을 구성하고, 전문가 검토를 통해 문항의 적절성을 확인한 뒤, 예비 조사를 거쳐 본 조사에 사용했다는 흐름을 제시한다.
연구 절차가 복잡한 경우에는 그림~\ref{fig:tikz-process-example}처럼 절차도나 흐름도를 함께 제시하면 좋다.

\begin{figure}[htbp]
    \centering
    \footnotesize
    \begin{tikzpicture}[
        process/.style={
            draw,
            rounded corners,
            minimum width=2.5cm,
            minimum height=0.9cm,
            align=center
        },
        arrow/.style={->, thick}
    ]
        \node[process] (review) {문헌 검토};
        \node[process, right=0.55cm of review] (design) {연구 설계};
        \node[process, right=0.55cm of design] (analysis) {자료 분석};
        \node[process, right=0.55cm of analysis] (conclusion) {결론 도출};

        \draw[arrow] (review) -- (design);
        \draw[arrow] (design) -- (analysis);
        \draw[arrow] (analysis) -- (conclusion);
    \end{tikzpicture}
    \caption{연구 절차도 예시}
    \label{fig:tikz-process-example}
\end{figure}

그림~\ref{fig:tikz-process-example}과 같은 절차도는 TikZ로 직접 그릴 수 있다. 작성 방법은 \href{../manual/7-KNUE_THESIS_TEMPLATE_USAGE.md}{manual/7 문서}의 "수식과 TikZ" 절을 참고한다.

\subsection{자료 수집}

% 연구에서 사용한 자료와 수집 방법을 설명합니다.
% 설문지, 검사 도구, 면담, 관찰, 수업 산출물, 로그 데이터, 공개 데이터,
% 문헌 자료 등을 구체적으로 밝힙니다.
% 검사 도구나 설문지를 사용한 경우 문항 수, 척도, 하위 영역, 실시 시간을 씁니다.

자료 수집 절에서는 과학교육 연구 결과를 도출하기 위해 어떤 자료를 어떻게 모았는지 설명한다.
개념 검사, 탐구 능력 검사, 과학 태도 설문지, 면담 자료, 수업 관찰 기록, 활동지와 보고서 같은 학습 산출물, 온라인 학습 로그, 공개 데이터, 문헌 자료 등 자료의 종류를 구체적으로 밝힌다.
자료 수집 시기, 학교와 교실 맥락, 수업 단원, 차시, 소요 시간, 실시 방법, 연구자와 교사의 역할도 필요한 범위에서 제시한다.

검사 도구나 설문지를 사용한 경우에는 문항 수, 응답 척도, 하위 영역, 점수 산출 방식, 실시 시간을 설명한다.
면담을 수행한 경우에는 면담 방식, 면담 횟수, 1회 면담 시간, 주요 질문 영역, 녹음과 전사 여부를 밝힌다.
자료 수집 방법이 충분히 구체적이어야 독자가 연구 결과의 신뢰성과 해석 가능성을 판단할 수 있다.

\subsection{자료 분석}

% 수집한 자료를 어떻게 분석했는지 설명합니다.
% 양적 자료라면 기술통계, t 검정, ANOVA, 상관분석, 회귀분석, 요인분석 등
% 사용한 통계 방법과 소프트웨어를 씁니다.
% 질적 자료라면 전사, 개방 코딩, 범주화, 주제 도출, 평정자 협의 절차 등을 씁니다.
% 분석 기준표나 코딩틀이 있다면 표로 제시하면 좋습니다.

자료 분석 절에서는 수집한 자료를 어떤 기준과 절차에 따라 분석했는지 설명한다.
양적 연구에서는 기술통계, 빈도 분석, t 검정, 분산분석, 상관분석, 회귀분석 등 사용한 통계 방법을 밝히고, 분석에 사용한 프로그램과 유의수준을 제시한다.
또한 결측치 처리, 이상치 확인, 정규성 검토처럼 분석 전에 수행한 자료 점검 절차가 있다면 함께 쓴다.

질적 연구에서는 녹음 자료의 전사, 반복 읽기, 의미 단위 추출, 개방 코딩, 범주화, 주제 도출의 과정을 설명한다.
여러 연구자가 함께 코딩했다면 코딩자 수, 의견 불일치 조정 방식, 평정자 간 일치도 또는 협의 절차를 제시한다.
혼합 연구에서는 양적 자료와 질적 자료를 각각 어떻게 분석했는지, 그리고 두 자료를 어느 단계에서 어떻게 통합했는지 밝혀야 한다.

\subsection{타당도와 신뢰도 확보}

% 연구 도구와 분석 결과의 타당도, 신뢰도를 어떻게 확보했는지 설명합니다.
% 예: 전문가 검토, 내용 타당도 지수, 예비 조사, Cronbach's alpha,
% 평정자 간 일치도, 삼각검증, 참여자 확인, 동료 검토 등

타당도와 신뢰도 확보 절에서는 연구 도구, 자료 수집 과정, 분석 결과가 믿을 만한지 확인하기 위해 어떤 조치를 했는지 설명한다.
양적 연구에서는 전문가 검토, 예비 조사, 문항 분석, 내용 타당도 지수, Cronbach's alpha와 같은 신뢰도 계수를 제시할 수 있다.
질적 연구에서는 삼각검증, 참여자 확인, 동료 검토, 연구자 성찰 기록, 풍부한 맥락 설명 등을 통해 해석의 신뢰성을 높일 수 있다.

이 절에서는 단순히 "타당도와 신뢰도를 확보하였다"라고 쓰지 말고, 누가 어떤 기준으로 검토했는지, 어떤 절차를 거쳤는지, 그 결과 연구 도구나 분석이 어떻게 수정되었는지를 구체적으로 밝힌다.

\subsection{연구 윤리}

% 연구 참여자의 권리와 개인정보를 어떻게 보호했는지 설명합니다.
% 연구 참여 동의, 익명화, 자료 보관 방법, 민감 정보 처리,
% IRB 승인 여부 또는 심의 면제 여부를 필요에 따라 씁니다.

연구 윤리 절에서는 연구 참여자의 권리와 개인정보를 어떻게 보호했는지 설명한다.
사람을 대상으로 한 연구에서는 연구 목적, 참여 절차, 예상 소요 시간, 참여자의 권리, 중도 철회 가능성을 안내하고 동의를 받았는지 밝힌다.
수집한 자료는 익명화 또는 가명 처리하고, 연구 목적 외에는 사용하지 않으며, 보관 기간과 폐기 방법을 명시한다.

기관생명윤리위원회(IRB) 심의가 필요한 연구라면 승인 여부와 승인 번호를 제시한다.
심의 면제 대상인 경우에도 왜 면제에 해당하는지, 연구 참여자에게 위험이 없도록 어떤 조치를 했는지 설명하는 것이 좋다.
특히 학생, 미성년자, 교사, 환자 등 취약할 수 있는 참여자가 포함되는 경우에는 동의 절차와 개인정보 보호를 더 구체적으로 적는다.

이 장에서 제시한 연구 설계, 자료 수집, 자료 분석 절차는 다음 장의 연구 결과를 해석하는 기준이 된다.
따라서 결과 장에서는 여기에서 제시한 연구 문제와 분석 절차의 순서에 맞추어 결과를 배열하면 문서 전체의 흐름이 자연스럽다.
